\documentclass[11pt]{article}
\usepackage[a4paper,margin=1in]{geometry}
\usepackage{amsmath,amssymb,booktabs,graphicx,hyperref,xcolor}
\title{Autonomous Diffusion: locked-test verification on CIFAR-10\\
       {\large 10{,}000 samples / cell, 5 seeds, 17 samplers, 510 runs}}
\author{thermau5}
\date{\today}
\begin{document}
\maketitle

\paragraph{Status.} \textbf{Locked test complete.} The contract
(\texttt{outputs/freeze\_record.yaml}, sha256 \texttt{450925ba0791\ldots})
was frozen on 2026-05-24 and not modified during the 45.77\,h evaluation.
The critic guards refused any drift; \(510/510\) runs succeeded, zero
failures. Per-cell numbers below are mean \(\pm\) SEM over 5 declared
seeds \(\{0, 1, 2, 3, 4\}\). No tuning post-freeze.

\section{Theory}
\label{sec:theory}

\paragraph{Setup.} Let \(p_0=\mathcal N(0,I)\) be the prior and
\(p_1 = p_{\mathrm{data}}\) the data distribution. The generative ODE
in noise coordinate \(\sigma\) is
\[
\frac{dx}{d\sigma} = g\bigl(x,\sigma;\theta\bigr),
\qquad
g(x,\sigma;\theta) = \frac{x - D_\theta(x,\sigma)}{\sigma},
\]
with \(D_\theta\) the pretrained EDM denoiser. A \emph{sampler} is the discretised controlled flow
\[
x_{i+1} = F_s\bigl(x_i,\sigma_i,\sigma_{i+1};\theta,\mathcal H_i\bigr),
\quad i = 0,\ldots,K-1,
\]
specified by the tuple
\[
u = (\theta,\ \pi,\ s,\ m,\ N)
\]
where \(\theta\) is the backbone, \(\pi\) is the path/coupling family, \(s\)
the solver core (update law \(F_s\)), \(m(\sigma)\) the sampling-grid
density inducing \(\sigma_0>\sigma_1>\cdots>\sigma_K=0\), and \(N\)
the cost budget (NFE / wall-clock).

\paragraph{Certified KL/NLL risk.} The certificate decomposes the
terminal-risk bound additively:
\[
R(u) \;\le\; Q_s(u) \;=\; Q_0
\;+\;
\frac{c_{\mathrm{stat}}}{n}\!\int_0^{B}\!\frac{a(r)}{\rho(r)}\,dr
\;+\;
c_{\mathrm{disc}}\!\int_0^{B}\!\frac{d_s(r)}{m(r)^p}\,dr,
\]
where \(\rho\) is the training/path-time allocation (frozen by the
pretrained net for this work), \(a(r)\) is the per-\(r\) statistical
weight, \(d_s(r)\) is the solver-specific discretisation density at noise
level \(r\), and \(p\) is the solver order (Heun: \(p=2\)).

\paragraph{Solver-specific discretisation density.} \(d_s(\sigma)\) is
the per-\(\sigma\) local-truncation-error density of the chosen solver
\(s\). Empirically (full recipe in Appendix~A): take a batch of
in-distribution noisy images at \(\sigma_a\), step them once with the
actual update law \(F_s\) to \(\sigma_b\), compare against a
16-substep Heun reference, and normalise by interval length:
\[
\boxed{\;
d_s(\sigma_i)
=
\frac{1}{|\sigma_a-\sigma_b|}\,
\frac{1}{B}\sum_{b=1}^{B}
\bigl\| F_s\bigl(x_b;\sigma_a\!\to\!\sigma_b\bigr)
       - \mathrm{Heun}^{(16)}\bigl(x_b;\sigma_a\!\to\!\sigma_b\bigr) \bigr\|_{\mathrm{RMS}}
\;}
\]
on a uniform-log-\(\sigma\) grid of \(M=32\) intervals over
\([\sigma_{\min},\sigma_{\max}]=[0.002,80]\). The reference is solver-
independent; the test step is the actual \(F_s\), so \(d_s\) is
solver-specific by construction.

\paragraph{Closed-form Pareto-optimal schedule.} Solving the certificate's
fixed-NFE constrained problem
\[
\min_{m}\!\int\!\frac{d_s(\sigma)}{m(\sigma)^p}\,d\sigma
\quad\mathrm{s.t.}\quad
\int m(\sigma)\,d\sigma = N
\]
via Lagrangian KKT gives the closed form:
\[
\boxed{\;\;
m^\star_s(\sigma)
\;\propto\;
d_s(\sigma)^{\,1/(p+1)}
\;\;}
\]
For each solver core \(s\in\{\mathrm{Heun},\mathrm{DPM\text{-}Solver\text{++}},\mathrm{UniPC},\mathrm{DEIS},\mathrm{Restart}\}\) we run \texttt{proposed\_<core>} with its actual update law and the certificate-optimal grid \(m^\star_s\) above. The implementation applies a Karras-style perceptual weight \(w(\sigma)=\sigma^{-k}\) with \(k=2\) (the only validation-tuned hyperparameter) so the effective form is \(m^\star_s(\sigma)\propto\!\bigl(d_s(\sigma)\!\cdot\!w(\sigma)\bigr)^{1/(p+1)}\).

\paragraph{Scope (Appendix A correction).} The boxed pointwise form
applies to single-step / pointwise-defect solvers (RK-class: Heun, EDM-Euler, DPM-Solver-2) and to multistep solvers whose leading
defect is empirically pointwise (DPM-Solver++ 2M, UniPC --- Theorem A).
For \(q\)-step Adams--Bashforth solvers (DEIS tAB-2, PNDM/PLMS-4) the
local residual is the history-coupled interpolation-kernel error
\(e_i^{\mathrm{AB}} \propto g^{(q)}(\xi)\,\prod_{\ell=0}^{q-1}(\sigma-\sigma_{i-\ell})\),
which does not factor into a scalar \(d_s(\sigma)\). The certificate is
then replaced by a coupled-sequence functional
\(D_s^{\mathrm{AB}}(\Pi_N)\) and the optimal grid is found by
sequence optimisation \(\Pi_{N,s}^\star = \arg\min_{\Pi_N} D_s^{\mathrm{AB}}(\Pi_N)\) (Theorem B,
implemented in this work).

\paragraph{Headline empirical result.} On the locked CIFAR-10 test split (10k samples / cell, 5 seeds, 510 runs, contract sha256 \texttt{450925ba0791\ldots}):
\[
\boxed{\;\;
\bigl(\,s_\star,\;m^\star_{s_\star}\bigr) \;=\; \bigl(\mathrm{UniPC},\;\mathrm{Ours}\bigr)
\quad\text{owns the entire cost-honest Pareto frontier at}\quad
\mathrm{actual\ NFE}\in[5,32]
\;\;}
\]
with 6 of 9 frontier points held by (Ours, \(s\)) variants and the largest
single-cell improvement \(-31.7\) FID at Heun NFE\,=\,9.

\section{Comparison Hierarchy}

A generative sampler is the tuple
\[
u = (\theta,\ \pi,\ s,\ m,\ N)
\]
where \(\theta\) is the backbone (architecture + weights), \(\pi\) is the path/coupling family (which generative ODE/SDE the backbone parameterises), \(s\) is the solver core (the discrete update law \(F_s\)), \(m\) is the schedule/grid (where on the noise axis \(F_s\) is evaluated), and \(N\) is the cost budget (actual NFE or wall-clock). Our method is not a new generative model; it is per-solver-core schedule control:
\[
m^\star_s(\sigma) \propto d_s(\sigma)^{1/(p+1)}.
\]

We organise experiments by which variable changes:

\begin{center}
\scriptsize
\begin{tabular}{c l l p{5.0cm} l}
\toprule
Level & Variable & Fixed objects & Question & Status \\
\midrule
1 & \(s\)      & \(\theta,\pi,m_s^{default}\)         & Strongest solver baseline? & \textbf{done} \\
2 & \(m\)      & \(\theta,\pi,s\)                     & Does \(m^\star_s\) beat \(m_s^{default}\)? & \textbf{done} \\
3 & \(\pi\)    & comparable backbone/compute          & Does path-straightening beat schedule control? & \emph{open} \\
4 & \(\theta\) & none / broad protocol                & Full generative-model leaderboard & \emph{out of scope} \\
\bottomrule
\end{tabular}
\end{center}

\noindent
The locked test in this report covers Level 1 + Level 2 under
\(\theta = \theta_{\mathrm{EDM}}\), \(\pi = \pi_{\mathrm{EDM}}\). It does
not include Level 3 (Rectified Flow / OT-CFM / MSFM / OFM), which is the
explicit next round before any path-straightening claim.

\section{Baseline audit}

For each sampler used in the locked test, the per-experiment fixed/variable
objects:

\begin{center}
\scriptsize
\begin{tabular}{l l l l l l l}
\toprule
sampler id & path \(\pi\) & backbone \(\theta\) & solver core \(s\) & baseline grid \(m_s^{default}\) & paired \((\)Ours, \(s)\) & actual NFE \\
\midrule
\texttt{edm\_euler}        & EDM diffusion & pretr.\ EDM & Euler         & Karras \(\rho{=}7\)       & --                    & \(K\) \\
\texttt{edm\_heun}         & EDM diffusion & pretr.\ EDM & Heun          & Karras \(\rho{=}7\)       & \texttt{proposed\_heun} & \(2K{-}1\) \\
\texttt{karras\_schedule}  & EDM diffusion & pretr.\ EDM & Heun          & Karras \(\rho{=}7\)       & \texttt{proposed\_heun} & \(2K{-}1\) \\
\texttt{uniform\_schedule} & EDM diffusion & pretr.\ EDM & Heun          & uniform-log               & \texttt{proposed\_heun} & \(2K{-}1\) \\
\texttt{ddim}              & EDM diffusion & pretr.\ EDM & DDIM (1st)    & DDPM \(\beta\)-derived    & --                    & \(K\) \\
\texttt{ddpm\_ancestral}   & EDM diffusion & pretr.\ EDM & DDPM ancestral& DDPM \(\beta\)-derived    & --                    & \(K\) \\
\texttt{dpm\_solver}       & EDM diffusion & pretr.\ EDM & DPM-Solver-2  & Karras \(\rho{=}7\)       & --                    & \(2K{-}1\) \\
\texttt{dpm\_solver\_pp}   & EDM diffusion & pretr.\ EDM & DPM-Solver++ 2M & Karras \(\rho{=}7\) in our impl. & \texttt{proposed\_dpmpp} & \(K\) \\
\texttt{unipc}             & EDM diffusion & pretr.\ EDM & UniPC 2PC     & Karras \(\rho{=}7\) in our impl. & \texttt{proposed\_unipc} & \(K\) \\
\texttt{deis}              & EDM diffusion & pretr.\ EDM & DEIS tAB-2    & Karras \(\rho{=}7\) in our impl. & \texttt{proposed\_deis} & \(K\) \\
\texttt{pndm}              & EDM diffusion & pretr.\ EDM & PNDM/PLMS-4   & Karras \(\rho{=}7\) in our impl. & --        & \(K{+}\text{warmup}\) \\
\texttt{restart}           & EDM diffusion & pretr.\ EDM & Restart wrap.\ Heun & Xu et al.\ \(\sigma\)-band & \texttt{proposed\_restart} & base + cycles \\
\bottomrule
\end{tabular}
\end{center}

\noindent
\textbf{Wording note.} The ``baseline grid'' column is the standard
Karras / DDPM-derived grid used in \emph{our EDM-compatible
implementation}, not necessarily the exact grid each method's original
paper or release configures. The paired \((\)Ours, \(s)\) variant runs
the same solver core on the certificate-optimal grid \(m_s^\star\) from
this work. ``--'' means we have no paired variant (DDIM, DDPM, EDM-Euler,
DPM-Solver-2, PNDM are baseline-only; the latter four are subsumed by
their multistep / corrector variants on the proposed side).

\section{Schedule (\(m\)) taxonomy}

The grid \(m(r)\) inducing \(\sigma_0 > \sigma_1 > \cdots > \sigma_K = 0\)
takes one of five classes. Issue \#2 in the repo demands this distinction
explicitly:

\begin{center}
\small
\begin{tabular}{l l p{6.5cm}}
\toprule
class & example & description \\
\midrule
fixed heuristic            & Karras \(\rho{=}7\), uniform-log-\(\sigma\)   & fixed before validation/test \\
tuned parametric           & Karras \(\rho^\star\)                         & small parameter vector optimised on validation \\
learned scheduler          & \(m_\phi\)                                    & trained schedule model/policy \\
adaptive scheduler         & \(m(x,r)\)                                    & sample/trajectory-dependent online schedule \\
\textbf{nonparametric-calibrated} & \textbf{ours \(m_s^\star\)}            & full calibrated profile \(d_s(r)\) + closed-form certificate optimiser \\
\bottomrule
\end{tabular}
\end{center}

\paragraph{Canonical statement (verbatim).}
Ours is a \textbf{nonparametric-calibrated scheduler}: it estimates a
solver-specific difficulty profile \(d_s(r)\) over the noise coordinate
and then computes the certificate-optimal grid
\(m_s^\star(r) \propto d_s(r)^{1/(p+1)}\). It is nonparametric \emph{relative
to low-dimensional schedule families} such as Karras \(\rho\), because the
calibrated object is a full profile \(d_s(r)\), not a small parameter
vector. It is calibrated \emph{rather than learned} because \(d_s(r)\) is
measured from local solver error against a reference solver, not trained
end-to-end by FID/test loss.

\paragraph{What the locked test already covers.} fixed-heuristic and
nonparametric-calibrated only:
\begin{itemize}
\item Karras \(\rho{=}7\) (\texttt{karras\_schedule}, \texttt{edm\_heun}, \texttt{edm\_euler}, \texttt{dpm\_solver\_pp}, \texttt{unipc}, \texttt{deis}, \texttt{pndm}, \texttt{restart} all run the Karras \(\rho{=}7\) grid in our EDM-compatible implementation);
\item Uniform-log-\(\sigma\) (\texttt{uniform\_schedule});
\item DDPM-\(\beta\)-derived (\texttt{ddim}, \texttt{ddpm\_ancestral});
\item Solver implementation defaults (\texttt{dpm\_solver}, \texttt{restart});
\item Nonparametric-calibrated \(m_s^\star\) (\texttt{proposed\_heun}, \texttt{proposed\_dpmpp}, \texttt{proposed\_unipc}, \texttt{proposed\_deis}, \texttt{proposed\_restart}).
\end{itemize}

\paragraph{What the locked test does not yet cover.}
\begin{itemize}
\item \textbf{Tuned parametric.} No validation-optimised Karras \(\rho^\star\) baseline was included. The locked test uses the fixed \(\rho{=}7\) member of the Karras family for every Karras-grid sampler. \emph{Round 5a (in progress) addresses this with a \(\rho \in \{1,2,3,5,9,12,15\}\) sweep on UniPC.}
\item \textbf{Learned scheduler.} No \(m_\phi\) baseline (e.g.\ neural timestep policy).
\item \textbf{Adaptive scheduler.} No sample-dependent \(m(x,r)\) baseline.
\end{itemize}

Until a tuned parametric and a learned/adaptive schedule baseline are
either implemented or marked infeasible with citation, the headline
\((s,\ m_s^\star)\) win is to be interpreted as a win against
\emph{fixed-heuristic} schedules only, not against tuned/learned/adaptive
ones.

\paragraph{Literature schedule-optimisation baseline (Issue \#2 \S3.3).}
Closest published schedule-optimisation methods and their fit for our
fixed-EDM benchmark:

\begin{itemize}
\item \textbf{AYS \emph{(Align Your Steps)}} (Sabour, Hayat, Garg et al., ICLR 2024): optimises a
  timestep grid by minimising KL between the discrete sampler trajectory
  and a continuous-SDE reference. Released schedules: SDXL and SD~1.5
  only. CIFAR-10/EDM schedules are not publicly released; reproducing
  them requires re-running the AYS optimisation on the EDM net
  ($\sim$1--2 days of careful implementation).
\item \textbf{DDSS \emph{(Watson et al.\ 2022, ``Learning fast samplers
  for diffusion models'')}}: differentiable sampler search over
  schedule parameters. Public code is partial; the paper publishes
  CIFAR-10 schedules in its appendix but as fitted parameter vectors,
  not as drop-in code.
\item \textbf{DPM-Solver-v3} (Lu, Zhou et al., NeurIPS 2023): bundles ``empirical
  model statistics'' solver corrections with an associated schedule. Not
  a pure schedule baseline (changing the schedule alone does not isolate
  the schedule contribution).
\item \textbf{GENIE} (Dockhorn et al.\ 2022): distills a 2nd-order solver
  with its schedule. Same bundling issue.
\end{itemize}

For this report, we mark the AYS / DDSS comparison \textbf{infeasible
within the current single-GPU budget} (the dominant fixed cost is the
re-derivation pipeline, not the eventual evaluation): faithful AYS
re-derivation on EDM CIFAR-10 is a $\sim$1--2 day item. We choose to
\textbf{defer to a follow-up round} on the grounds that (a) the tuned
Karras-$\rho^\star$ baseline (\S\ref{sec:tuned-rho}) below is the
\emph{within-family} schedule comparison the issue most directly asks
for; (b) the bundled methods (DPM-Solver-v3, GENIE) cannot be used as
pure schedule baselines; (c) AYS/DDSS implementation would not change
the tuned-$\rho^\star$ comparison result we already get, only add one
row to the Level-2 schedule table.

\section{Level-2: paired same-solver wins}

For each solver core \(s\), the locked-test FID delta at matched actual NFE,
\(\Delta_s(K) = \mathrm{FID}(s, m_s^\star, K) - \mathrm{FID}(s, m_s^{default}, K)\),
negative \(\Rightarrow\) the certificate grid wins for this solver core
\emph{at fixed everything else}:

\begin{center}
\small
\begin{tabular}{l c c c c}
\toprule
solver core \(s\) & NFE\,9 & NFE\,15 & NFE\,23 & NFE\,35 \\
\midrule
Heun (vs Karras)         & \(-31.70\) & \( -1.86\) & \(+0.07\) & \(+0.07\) \\
\midrule
solver core \(s\) & NFE\,5 & NFE\,8 & NFE\,12 & NFE\,18 \\
\midrule
DPM-Solver++         & \(-16.95\) & \( -3.71\) & \(-0.75\) & \(-0.24\) \\
UniPC                & \(-18.88\) & \( -3.33\) & \(-0.47\) & \(-0.13\) \\
DEIS                 & \( -6.51\) & \( -1.27\) & \(+0.10\) & \(-0.01\) \\
\bottomrule
\end{tabular}
\end{center}

\noindent
\textbf{What this isolates.} Fix \(\theta = \theta_{\mathrm{EDM}}\),
\(\pi = \pi_{\mathrm{EDM}}\), and the solver core \(s\). Only the
grid \(m\) changes from \(m_s^{default}\) (Karras-style) to \(m_s^\star\)
(certificate-optimal). For 3 of 4 cores (Heun, DPM-Solver++, UniPC) the
certificate grid wins at every tested NFE from 5/9 up to 12. The DEIS
column shows the same low-NFE win because \texttt{proposed\_deis} uses
the shared \(d_{\mathrm{Heun}}\) grid (Theorem A); the corrected appendix's
sequence-level Theorem B is implemented but does not yet outperform the
shared heuristic empirically (pending a better terminal-amplification
estimator).
SEMs are \(\sim 0.025\) so every \(|\Delta| \ge 0.1\) is \(>4\sigma\)
separable.

\noindent
\textbf{What this does not test.} Whether changing \(\pi\) (the path /
coupling) on a different backbone wins more than changing \(m\) on the
fixed EDM backbone. That is Level 3 and is open.

\section{Headline (locked test)}

\begin{figure}[h]
\centering
\includegraphics[width=0.88\linewidth]{pareto_locked_zoom.png}
\caption{Locked-test cost-honest Pareto frontier on CIFAR-10, low-NFE region (actual NFE \(\le 35\)). \textbf{(Ours, UniPC)} owns the entire NFE=5 through NFE=32 corner; (Ours, Restart) takes NFE=33; baseline Restart picks up NFE \(\ge 39\).}
\end{figure}

\paragraph{Pareto frontier.}
\begin{center}
\small
\begin{tabular}{r l r r r}
\toprule
actual NFE & sampler & Clean-FID & SEM & owner \\
\midrule
  5 & (Ours, UniPC)    & 21.459 & 0.136 & \textbf{ours} \\
  8 & (Ours, UniPC)    &  9.183 & 0.070 & \textbf{ours} \\
 12 & (Ours, UniPC)    &  5.578 & 0.021 & \textbf{ours} \\
 18 & (Ours, UniPC)    &  4.664 & 0.022 & \textbf{ours} \\
 27 & PNDM             &  4.577 & 0.023 & baseline \\
 31 & Restart          &  4.486 & 0.030 & baseline \\
 32 & (Ours, UniPC)    &  4.447 & 0.024 & \textbf{ours} \\
 33 & (Ours, Restart)  &  4.387 & 0.021 & \textbf{ours} \\
 39 & Restart          &  4.339 & 0.026 & baseline \\
\bottomrule
\end{tabular}
\end{center}

\noindent
\textbf{(Ours)} variants occupy \textbf{6 of 9} Pareto-frontier points.
\textbf{(Ours, UniPC)} alone owns the entire NFE=5\,\(\to\)\,32 corner.

\paragraph{Pareto AUC.}
\[
\mathrm{AUC}_{\mathrm{locked}}\;(\log_{10}\mathrm{NFE} \in [4, 200]) \;=\; 10.755.
\]

% (Per-core within-solver-core wins table moved to Section 4
%  'Level-2: paired same-solver wins' added per Issue #1.)

\section{Per-sampler best Clean-FID}

\begin{center}
\small
\begin{tabular}{l r r r}
\toprule
sampler & best Clean-FID & SEM & at actual NFE \\
\midrule
Restart                & 4.339 & 0.026 &  39 \\
(Ours, Restart)        & 4.360 & 0.025 &  41 \\
EDM-Heun               & 4.398 & 0.023 &  65 \\
Karras schedule        & 4.398 & 0.023 &  63 \\
(Ours, Heun)           & 4.403 & 0.023 & 127 \\
Uniform-log schedule   & 4.403 & 0.024 &  63 \\
PNDM                   & 4.409 & 0.024 &  73 \\
\textbf{(Ours, UniPC)} & \textbf{4.414} & 0.023 & \textbf{64} \\
(Ours, DPM-Solver++)   & 4.420 & 0.023 &  64 \\
UniPC                  & 4.432 & 0.024 &  64 \\
DPM-Solver             & 4.436 & 0.023 &  65 \\
DPM-Solver++           & 4.450 & 0.023 &  64 \\
(Ours, DEIS)           & 4.463 & 0.023 &  64 \\
DEIS                   & 4.482 & 0.024 &  64 \\
EDM-Euler              & 5.135 & 0.029 &  64 \\
DDIM                   & 7.111 & 0.033 &  64 \\
DDPM ancestral         &13.292 & 0.044 &  64 \\
\bottomrule
\end{tabular}
\end{center}

\paragraph{Note on the high-NFE corner.} \texttt{Restart} achieves the
absolute best FID (4.339) at NFE=39 by paying for an extra noise-injection
cycle on top of an ODE solver; \texttt{proposed\_restart} mirrors that
structure but uses our certificate's grid for the base trajectory and lands
at 4.360 (within \(1\sigma\)). The certificate is competitive with the
Restart wrapper but doesn't beat it; lifting the certificate to also
optimize the wrapper-cycle structure is future work.

\section{Full Pareto figure}

\begin{figure}[h]
\centering
\includegraphics[width=0.88\linewidth]{pareto_locked.png}
\caption{Locked-test Pareto frontier across all 17 samplers and all NFE budgets. Log-log axes. The flat floor near FID \(=4.4\) at NFE \(\ge 32\) is the data-distribution / Clean-FID noise floor at 10k samples; tightening it would require 50k samples (~15\,GPU-days on a single 4090).}
\end{figure}

\section{Caveats and follow-ups}

\begin{itemize}
\item \textbf{10k vs 50k.} The locked-test sample count was reduced from
the original contract's 50{,}000 to 10{,}000 to fit a single-GPU budget
(documented in YAML; the relative ordering of samplers is preserved at 10k
since the saturated floor is well above SEM). A gold-standard 50k re-run
of the Pareto-frontier cells alone (\(\sim 9\) cells, 5 seeds) is a
2-3\,day follow-up.
\item \textbf{DEIS Theorem-B implementation (Round 3c).} The corrected PDF 1 appendix introduces sequence-level grid optimization \(\Pi_{N,s}^\star = \arg\min_{\Pi_N} D_s^{\mathrm{AB}}(\Pi_N)\) for AB-class solvers. This is implemented for DEIS tAB-2 (1k smoke: 48 / 38 / 35 / 33 at NFE 5/8/12/18). With the current \(g''(\sigma)\) finite-difference estimator and an \(A_i = \sigma^{-k}\) amplification heuristic, the sequence-\(\Pi^\star\) grid minimizes the right symbolic objective but is empirically worse than the shared \(d_{\mathrm{Heun}}\) heuristic (Round 3b 1k smoke: 40 / 23 / 19 / 17). A proper terminal-amplification estimator \(A_i\) is open. \texttt{proposed\_deis} therefore defaults to the shared \(d_{\mathrm{Heun}}\) grid, which is what the locked test reports.
\item \textbf{Wasserstein variant.} The Wasserstein OT certificate (PDF 2) optimising \(\nu\) is not yet implemented; the present result tests only the KL/NLL variant on \(m\).
\end{itemize}

\section{Conclusion}

\[
\boxed{
m^\star_s(\sigma) \propto d_s(\sigma)^{1/(p+1)} \text{ creates new Pareto-frontier points at low NFE on the locked test split.}
}
\]

The strongest single-sampler result is \((s = \mathrm{UniPC}, m = m^\star_s)\)
holding the entire NFE\,\(\in [5, 32]\) corner of the cost-honest Pareto
frontier on CIFAR-10. Six of nine frontier points are (Ours) variants.
The certificate's per-core lifting is empirically verified.

\end{document}
